\documentclass[aspectratio=169,spanish,xcolor=table]{beamer}
\usepackage{iftex}
\ifPDFTeX
  \usepackage[utf8]{inputenc}
  \usepackage[T1]{fontenc}
\fi
\usepackage[spanish,es-tabla,es-noshorthands]{babel}
\usepackage{helvet}
\renewcommand{\familydefault}{\sfdefault}
\usepackage{tikz,pgfplots,array,booktabs,ragged2e}
\pgfplotsset{compat=1.18}
\usetikzlibrary{arrows.meta,positioning,calc,shapes.geometric,fit}

\definecolor{iiiepeNavy}{HTML}{163B5C}
\definecolor{iiiepeBlue}{HTML}{1F6F9F}
\definecolor{iiiepeAqua}{HTML}{20B7A8}
\definecolor{iiiepeOrange}{HTML}{F28C28}
\definecolor{iiiepePaper}{HTML}{F6F8F7}
\definecolor{iiiepeInk}{HTML}{1E2A32}
\definecolor{iiiepePale}{HTML}{EEF5F8}
\definecolor{iiiepeMint}{HTML}{E7F6F4}
\definecolor{iiiepeWarm}{HTML}{FFF3E6}

\usetheme{default}
\usecolortheme{default}
\setbeamertemplate{navigation symbols}{}
\setbeamercolor{normal text}{fg=iiiepeInk,bg=white}
\setbeamercolor{structure}{fg=iiiepeNavy}
\setbeamercolor{frametitle}{fg=white,bg=iiiepeNavy}
\setbeamercolor{block title}{fg=white,bg=iiiepeBlue}
\setbeamercolor{block body}{fg=iiiepeInk,bg=iiiepePale}
\setbeamercolor{item}{fg=iiiepeOrange}

\setbeamertemplate{frametitle}{%
  \nointerlineskip%
  \begin{beamercolorbox}[wd=\paperwidth,ht=0.95cm,dp=0.22cm,leftskip=0.45cm,rightskip=0.45cm]{frametitle}%
    \begin{minipage}[c]{13.2cm}\usebeamerfont{frametitle}\large\bfseries\insertframetitle\end{minipage}%
    \hfill\begin{minipage}[c]{1.6cm}\raggedleft\includegraphics[height=0.56cm]{img/iiiepe/logo-iiiepe-blanco.png}\end{minipage}%
  \end{beamercolorbox}%
}

\setbeamertemplate{footline}{%
  \leavevmode%
  \hbox{%
    \begin{beamercolorbox}[wd=.78\paperwidth,ht=0.5cm,dp=0.18cm,leftskip=0.45cm]{frametitle}%
      \scriptsize\color{white}M. de la Cruz Muñoz · MAEM · Matrícula 00874 · SD26 · Actividad 7 · Modelo Bioecológico%
    \end{beamercolorbox}%
    \begin{beamercolorbox}[wd=.22\paperwidth,ht=0.5cm,dp=0.18cm,rightskip=0.45cm]{frametitle}%
      \hfill\scriptsize\color{white}\insertframenumber{} / \inserttotalframenumber%
    \end{beamercolorbox}%
  }%
  \vskip0pt%
}

\newcommand{\Pill}[2]{%
  \tikz[baseline=(x.base)]{\node[rounded corners=3pt,fill=#1!14,draw=#1!75!black,thick,inner xsep=6pt,inner ysep=4pt,align=center] (x) {\scriptsize\bfseries #2};}%
}

\title{El Modelo Bioecológico en Matemáticas}
\subtitle{Actividad 7 · Bloque III · Módulo 51813}
\author{Mtro. Martín Jonathan de la Cruz Muñoz · Matrícula 00874}
\institute{IIIEPE · Maestría en Aprendizaje y Enseñanza de las Matemáticas · Grupo A}
\date{19 de noviembre de 2026}

\begin{document}

\begin{frame}[plain]
  \begin{tikzpicture}[remember picture,overlay]
    \fill[iiiepeNavy] (current page.north west) rectangle ([xshift=10.3cm]current page.south west);
    \fill[iiiepeBlue] ([xshift=10.3cm]current page.north west) rectangle (current page.south east);
    \fill[iiiepeOrange] ([xshift=0.8cm,yshift=-6.45cm]current page.north west) rectangle ([xshift=9.3cm,yshift=-6.53cm]current page.north west);
    \node[anchor=north west,text width=8.7cm,align=left] at ([xshift=0.8cm,yshift=-1.15cm]current page.north west) {%
      {\color{iiiepeAqua!45!white}\small\bfseries MAESTRÍA EN APRENDIZAJE Y ENSEÑANZA DE LAS MATEMÁTICAS · SD26\par}%
      \vspace{0.35cm}
      {\color{white}\Large\bfseries \inserttitle\par}%
      \vspace{0.25cm}
      {\color{white!90}\normalsize \insertsubtitle\par}%
    };
    \node[anchor=south west,text width=8.7cm,align=left] at ([xshift=0.8cm,yshift=0.75cm]current page.south west) {%
      {\color{white}\normalsize\bfseries Mtro. Martín Jonathan de la Cruz Muñoz\par}%
      {\color{white!85}\scriptsize Matrícula 00874 · martin.delacruz@campus.iiiepe.edu.mx\par}%
      {\color{iiiepeAqua!35!white}\scriptsize Monterrey, Nuevo León · 19 de noviembre de 2026\par}%
    };
    \node[anchor=center,text width=5.0cm,align=center] at ([xshift=13.15cm,yshift=-4.5cm]current page.north west) {%
      \includegraphics[width=3.6cm]{img/iiiepe/logo-iiiepe-blanco.png}\par\vspace{0.45cm}%
      {\color{white}\large\bfseries IIIEPE\par}%
      \vspace{0.1cm}%
      {\color{white!85}\scriptsize Fundamentos para la enseñanza\\y el aprendizaje I\par}%
    };
  \end{tikzpicture}
\end{frame}
\begin{frame}{Modelo Bioecológico del Desarrollo Humano (PPCT)}
  \begin{block}{Tesis de Bronfenbrenner y Morris (2006)}
    El aprendizaje matemático emerge de interacciones continuas, recíprocas y progresivamente complejas entre una persona activa y contextos ecológicos anidados a través del tiempo.
  \end{block}
  \vspace{0.3cm}
  \centering
  \Pill{iiiepeNavy}{P · Proceso proximal} \hspace{0.2cm}
  \Pill{iiiepeBlue}{P · Persona} \hspace{0.2cm}
  \Pill{iiiepeAqua!75!black}{C · Contexto} \hspace{0.2cm}
  \Pill{iiiepeOrange!85!black}{T · Tiempo}
\end{frame}

\begin{frame}{1--2. Interacción PPCT y Crítica al Determinismo Individual}
  \begin{columns}[T]
    \column{0.48\textwidth}
      \begin{block}{1. Persona en contexto y experiencia vivida}
        \scriptsize
        \begin{itemize}
          \item La persona aporta disposiciones generadoras/disruptivas, recursos cognitivos y demandas sociales.
          \item Cuenta tanto el entorno objetivo como su vivencia fenomenológica (\textit{experiencing}).
        \end{itemize}
      \end{block}
    \column{0.48\textwidth}
      \begin{block}{2. Crítica al mérito aislado (Bronfenbrenner y Ceci)}
        \scriptsize
        \begin{itemize}
          \item Los potenciales no se actualizan sin procesos proximales sostenidos.
          \item Atribuir el fracaso únicamente a «capacidad o esfuerzo» oculta desigualdades de contexto y andamiaje.
        \end{itemize}
      \end{block}
  \end{columns}
\end{frame}

\begin{frame}{3--4. Comprender Contextos Reales en la Práctica Docente}
  \begin{columns}[T]
    \column{0.48\textwidth}
      \textbf{\color{iiiepeNavy}3. ¿Por qué importa al docente?}
      \begin{itemize}\scriptsize
        \item Evita etiquetar cansancio, trabajo o falta de equipo como «apatía».
        \item Orienta apoyos escolares realistas para proteger el derecho a aprender \textbf{(Tong y An, 2023)}.
      \end{itemize}
    \column{0.48\textwidth}
      \textbf{\color{iiiepeBlue}4. Mi contexto en CONALEP/CECyTE e Ingeniería:}
      \begin{itemize}\scriptsize
        \item Estudiantes con traslados largos, empleo parcial y acceso dispar a cómputo.
        \item Decisión: guías accesibles, contextos laborales auténticos y colaboración dentro del horario de clase.
      \end{itemize}
  \end{columns}
\end{frame}

\begin{frame}{5. Interacciones de Aula y Procesos Proximales (Proposición I)}
  \begin{block}{Condiciones para que la clase genere desarrollo matemático}
    \begin{itemize}\small
      \item \textbf{Actividad activa:} El estudiante manipula símbolos, datos y problemas; no solo escucha.
      \item \textbf{Regularidad y complejidad creciente:} Secuencias continuas que avanzan más allá de la repetición mecánica.
      \item \textbf{Reciprocidad bidireccional:} El docente escucha conjeturas y ajusta devoluciones en diálogo abierto.
      \item \textbf{Interacción con objetos y símbolos:} Instrumentos, gráficas y modelos físicos reales.
    \end{itemize}
  \end{block}
\end{frame}

\begin{frame}{6. Sistemas Contextuales Concéntricos}
  \centering
  \begin{tikzpicture}[font=\tiny, scale=0.76, transform shape]
    \draw[draw=iiiepeNavy,fill=iiiepePale,thick] (0,0) ellipse (6.8cm and 4.1cm);
    \draw[draw=iiiepeBlue,fill=white,thick] (0,0) ellipse (5.3cm and 3.2cm);
    \draw[draw=iiiepeAqua!80!black,fill=iiiepeMint,thick] (0,0) ellipse (3.8cm and 2.3cm);
    \draw[draw=iiiepeOrange!90!black,fill=iiiepeWarm,thick] (0,0) ellipse (2.3cm and 1.4cm);
    \node[fill=iiiepeNavy,text=white,rounded corners=2pt,text width=2.8cm,align=center,inner sep=3pt] at (0,-0.25)
      {\textbf{PROCESOS PROXIMALES}\\Persona $\leftrightarrow$ Símbolos y Pares};
    \node[fill=iiiepeOrange!90!black,text=white,rounded corners=2pt,font=\tiny\bfseries,inner sep=2pt] at (0,1.05) {Microsistema: Aula, taller, hogar};
    \node[fill=iiiepeAqua!80!black,text=white,rounded corners=2pt,font=\tiny\bfseries,inner sep=2pt] at (0,1.92) {Mesosistema: Relación escuela--hogar--taller};
    \node[fill=iiiepeBlue,text=white,rounded corners=2pt,font=\tiny\bfseries,inner sep=2pt] at (0,2.80) {Exosistema: Horarios, presupuesto, empleo};
    \node[fill=iiiepeNavy,text=white,rounded corners=2pt,font=\tiny\bfseries,inner sep=2pt] at (0,3.68) {Macrosistema: Cultura matemática, NEM, equidad};
  \end{tikzpicture}
\end{frame}

\begin{frame}{Factores Concretos por Sistema en Matemáticas}
  \scriptsize\renewcommand{\arraystretch}{1.35}
  \begin{tabular}{p{2.3cm}p{5.5cm}p{5.2cm}}
    \toprule
    \rowcolor{iiiepePale} \textbf{Sistema} & \textbf{Ejemplo situado en plantel técnico} & \textbf{Impacto didáctico} \\
    \midrule
    Microsistema & Clima de respeto ante el error y diálogo en binas. & Activa participación y seguridad matemática. \\
    Mesosistema & Vínculo entre matemáticas y taller de mecatrónica. & Da sentido y coherencia interdisciplinar. \\
    Exosistema & Disponibilidad de instrumentos y horarios laborales. & Requiere alternativas accesibles en clase. \\
    Macrosistema & Mitos sobre «talento innato» y políticas NEM. & Exige enfoque inclusivo contra exclusión. \\
    \bottomrule
  \end{tabular}
\end{frame}

\begin{frame}{Cronosistema: Microtiempo, Mesotiempo y Macrotiempo}
  \begin{itemize}\normalsize
    \item \textbf{\color{iiiepeNavy}Microtiempo:} Continuidad dialógica dentro de cada secuencia de clase sin saltos bruscos ni interrupciones punitivas.
    \item \textbf{\color{iiiepeBlue}Mesotiempo:} Seguimiento sistemático durante semanas del semestre para consolidar abstracción y modelación.
    \item \textbf{\color{iiiepeAqua!75!black}Macrotiempo:} Transformaciones tecnológicas e industriales que exigen habilidades analíticas y adaptativas nuevas.
  \end{itemize}
\end{frame}

\begin{frame}{Decisiones Ecológicas para Transformar el Aula}
  \begin{itemize}\normalsize
    \item Cuidar la \textbf{reciprocidad} escuchando cómo justifica cada equipo su modelo lineal o medición.
    \item Tejer el \textbf{mesosistema} vinculando ejercicios matemáticos con prácticas de taller e ingeniería.
    \item Neutralizar barreras del \textbf{exosistema} evitando depender de dispositivos privados inaccesibles fuera del horario escolar.
  \end{itemize}
\end{frame}

\begin{frame}{Conclusiones de la Actividad 7}
  \begin{itemize}\normalsize
    \item El modelo PPCT explica rigurosamente cómo persona, procesos proximales, contextos y tiempo configuran el aprendizaje matemático.
    \item Cuestionar el determinismo individual compromete a diseñar entornos ecológicos equitativos.
    \item Este marco se integra directamente en el proyecto final de la Actividad 8.
  \end{itemize}
\end{frame}

\begin{frame}{Referencias Principales (APA 7)}
  \scriptsize\setlength{\parskip}{5pt}
  \begin{itemize}
    \item Bronfenbrenner, U. (1987). \textit{La ecología del desarrollo humano}. Paidós.
    \item Bronfenbrenner, U., y Morris, P. A. (2006). The bioecological model of human development. En \textit{Handbook of child psychology} (pp. 793--828). Wiley.
    \item Collin, J., y Quigley, A. (2021). \textit{Teacher feedback to improve pupil learning}. EEF.
    \item CONALEP. (2025). \textit{Perfil de egreso: Profesional Técnico-Bachiller en Mecatrónica}. Secretaría Académica.
    \item SEP. (2023). \textit{Programa de estudios: Pensamiento Matemático I}. SEMS.
    \item Tong, X., y An, X. (2023). A review of empirical studies applying Bronfenbrenner's bioecological theory. \textit{Frontiers in Psychology}, 14, 1233925.
    \item UNESCO. (2015). \textit{Replantear la educación: ¿Hacia un bien común mundial?} UNESCO.
  \end{itemize}
\end{frame}

\end{document}
